% META: {"id": "1NSI-AGT-EX-002", "chapitre": "1NSI-ALGO-PARCOURS-TRIS", "type_objet": "exercice", "capacites": ["P-ALGO-01A"], "parcours": 1, "competences": ["programmer", "analyser"], "duree_min": 10, "mode_creation": "ex_nihilo", "sources_inspiration": [], "fichier_tex": "chapitres/1NSI-ALGO-PARCOURS-TRIS/exercices/1NSI-AGT-EX-002.tex", "corrige_tex": "chapitres/1NSI-ALGO-PARCOURS-TRIS/corriges/1NSI-AGT-CO-002.tex", "status": "verified"}
% BEGIN-VERIFY
% def recherche_toutes_occurrences(tableau, valeur):
%     indices = []
%     for i in range(len(tableau)):
%         if tableau[i] == valeur:
%             indices.append(i)
%     return indices
% assert recherche_toutes_occurrences([3, 7, 3, 3, 9], 3) == [0, 2, 3]
% assert recherche_toutes_occurrences([3, 7, 3, 3, 9], 5) == []
% END-VERIFY
\begin{exercice}{1NSI-AGT-EX-002}{1}{10}
La fonction \lstinline{recherche_occurrence} du cours ne renvoie que le premier indice
trouvé.
\begin{enumerate}
  \item Écrire une fonction \lstinline{recherche_toutes_occurrences(tableau, valeur)} qui renvoie la liste de \textbf{tous} les indices où \lstinline{valeur} apparaît.
  \item Tester cette fonction sur \lstinline{[3, 7, 3, 3, 9]} avec \lstinline{valeur=3}, puis avec \lstinline{valeur=5}.
  \item Quel est le coût, dans le pire cas, de cette fonction (en fonction de la taille $n$ du tableau) ? Est-il différent de celui de \lstinline{recherche_occurrence} ? Justifier.
\end{enumerate}
\end{exercice}
